%%%%%%%%%%%%%%%%%%%%%%%%%%%%%%%%%%%%%%%%%%%%%%%%%%%%%%%%%%%%%%%%%%%%%%%%%%%%%%%%%%%%%%
\section{The \system}
\label{sec:log}
%%%%%%%%%%%%%%%%%%%%%%%%%%%%%%%%%%%%%%%%%%%%%%%%%%%%%%%%%%%%%%%%%%%%%%%%%%%%%%%%%%%%%%
\system exposes to validators a Byzantine fault-tolerant log abstraction built around the \primitive primitive.
Validators can request to add blocks to the log.
Every validator 
can read from the log, or subscribe to asynchronous notifications when slots are finalized. 

Eventually, every slot is \emph{finalized} with some value.
Values can be blocks, or a special $\bot$ value.
Each slot has two finalization events, \commit and \eagercommit. 

More specifically, we define the log interface as follows:
\begin{itemize}
	\item $\init()$: initializes the log.
	\item $\propose(b)$: proposes to add block $b$.
	\item $\eagercommit(b)$: finalizes $b$.
	\item $\commit(b)$: commits $b$, its non-committed, slotless causal predecessors, and all blocks in its preceding slots in the log.
\end{itemize}

To propose a block $b$, a validator $p$ invokes $\bbcacast(id,sn,b)$ for a new instance $id$ and the next slot $sn$ available for $p$.
Policies for slot allocation to validators are discussed in \Cref{sec:ticketmaster}.


A \emph{block} becomes finalized when all the blocks in its causal references are finalized and either (1) is \primitivesmall-delivered, with or without a slot number, or (2) is decided by the fallback consensus protocol as the value for some slot.
A \emph{slot} $sn$ becomes finalized when a block becomes finalized for $sn$ or a $\bot$ value is decided by the fallback consensus protocol for $sn$. 
The slot becomes committed when the slot is finalized and its previous slot is committed. 
The log grows infinitely and starts with a known committed block (``genesis'').

When a slot becomes finalized at a validator,
\system triggers a \eagercommit notification to the validator. 
The finalized slot value may contain a block, or $\bot$ to mark a ``hole''.
Subsequently, when the slot becomes committed, 
\system triggers a \commit notification.


\paragraph{Fallback finalization protocol.}
Recall that \primitive only guarantees a weak validity property for slot numbers (see \primitiveproperty6 in \Cref{sec:bbca}).
To ensure that any slot is eventually finalized, 
and therefore all blocks added to the log can eventually become committed,
\system employs on a fallback finalization protocol which implements a consensus decision on stalled slots.

A validator $p$ maintains a timer for the first non-finalized slot it knows of as $\snmin$.
Upon finalizing $\snmin$, $p$ restarts a timer for the next non-finalized position.
If the timer expires, $p$ invokes the fallback finalization protocol by triggering the $\bbcaabort$ event for $\snmin$
and subscribes for $\bbcaannounce$ on $\snmin$.

Upon a $\bbcaannounce(sn, b, C_R,\mathcal{Q})$ event, $p$ \primitive's a block which includes in its metadata a $\finalize$ flag and the information $sn, b, C_R$.
If $p$ has not sent a \ready message for $sn$, then $C_R=\emptyset$ and $b=\bot$.

\Cref{alg:bbca-log} describes the log protocol using a consensus protocol with an external validity predicate (\vbc) as a ``blackbox''. 
\vbc exposes two events, $\bcpropose(v,Q)$ is triggered by a validator to participate in the protocol with a value $v$, and $\bcdecide(v)$.
We assume that $\vbc$ satisfies the following properties: if all correct validators participate with some value, eventually, all validators decide a value $v$, which is the value with which some validator participated (liveness) and, no two validators decide differently (safety).
Moreover, any value decided by a correct validator satisfies $\mathcal{Q}$.

Once a validator has \primitivesmall-delivered a set of $2f+1$ 
blocks with a $\finalize$ flag for some sequence number $sn$, it invokes $\bcpropose([sn,b,\mathcal{C}], \mathcal{Q})$, where $\mathcal{C}$ the corresponding set of $2f+1$ ready certificates.
The proposal entails a non-$\bot$ block $b$ if any of the $2f+1$ ready certificate is non-empty.
The external predicate $\mathcal{Q}$ enforces this condition to the value that correct validators deliver via $\vbtob$.
A correct validator finalizes slot $sn$ the first time it receives an event $\bcdecide([sn,b,\mathcal{C}])$ with some block $b$, as soon as it finalizes all $b'$ causal predecessors.

\paragraph{Fallback protocol implementation.}
There exist several consensus protocols in literature which satisfies the properties described above, such as HotStuff~\cite{yin2019hotstuff} and Tendermint~\cite{buchman2018latest}.
We implement the fallback protocol using the consensus framework of Fino~\cite{fino} on top of the DAG which is already populated by including causal references in \primitivesmall-delivered blocks, so that no additional messages are required to be exchanged.
Information pertaining to the consensus protocol is embedded in the metadata field of blocks.
In this way, every \primitive broadcast carries useful payload and allows the system to continue utilizing its bandwidth for transaction dissemination even when resolving non-finalized slots.

Briefly, the Fino consensus framework is structured in views and each view has a designated leader.
Upon entering a new view $r$, the leader broadcasts a proposal for view $r$ with \primitive, by including $proposal(r)$ in the metadata field of the block.
The proposal for view $r$ is justified if it has in its causal predecessors either $f+1$ votes for proposal $r-1$, or $2f+1$ complaints for round $r-1$.
When a proposal is added to the DAG, if it is justified then correct validators vote for it.
To vote for a proposal for view $r$, a validator includes $vote(r)$ in the metadata field of the block.
A consensus proposal is \emph{decided} in the DAG once it has $f+1$ votes succeeding it.
If the proposal for view $r$ is not decided in time, a correct validator \bbcacast\ a complaint by adding $complaint(r)$ in the metadata of the block.
Fino requires that messages by each validator are numbered and authenticated. If a validator delivers some message as the $k^{th}$ from a particular sender, then all other validators deliver the same message as the sender's $k^{th}$ message.
We use the external validity predicate $\mathcal{P}$ of \primitive to enforce this condition.

The first consensus proposal which follows $2f+1$ predecessor blocks with a $\finalize$ flag for some slot number $sn$, constitutes a consensus proposal to stop waiting for \primitive to finalize the slot $sn$.
It becomes decided by the Fino consensus protocol as described above.
Note that Fino proposals do not need to explicitly carry the set not ready certificates $\mathcal{C}$, as the reference the corresponding blocks which carry them as causal predecessors.
\Cref{fig:log-fallback} shows an example of a DAG instance for finalizing an expired slot number.


\begin{figure}[h]
  \centering
  \includegraphics[width=\columnwidth]{figures/DAG-consensus}
  \caption{Finalizing slot $4$: Blocks $5$,$7$, and $8$ indicate that validators the timers of $p0$, $p1$, and $p2$ expired for slot $4$.
  	Validator $p1$, who is a consensus leader for view $r$, implicitly proposes block $10$ for slot $4$. Blocks $12$ and $13$ are the votes for the leader proposal. }
  \label{fig:log-fallback}
\end{figure}

\paragraph{Causality.}
When validator $p$ creates a new block, 
it includes in it references to a set of \emph{causal predecessors}, a causal-frontier of previously \primitivesmall-delivered blocks which $p$ has not previously referenced in other blocks.
A reference to block $b$ is simply a pair $(sn,H(b))$ , where $H$ a collision resistant hash function and $sn$ the slot number with which $b$ was delivered, or $\bot$ if $b$ is a slotless block.
A validator before finalizing a block $b$ which it \primitivesmall-delivers, it finalizes all the blocks in 
$b$'s causal predecessors.
In  other words, finalizations are causal.
To ensure causal finalizations, we want to ensure that no correct validator delivers $b$ without delivering $b$'s causal predecessors first.
For this we use the external validity predicate of \primitive.
Moreover, a correct validator, upon receiving a \bbcacast event for a block $b$, starts a timer for all slot numbers of blocks in its $b$'s causal predecessors that are not finalized yet.
If the timer times out and the slot number is still not finalized, $p$ invokes $\bbcaabort(sn)$.
This timeout along with \primitive and the fallback finalization protocol ensure that if a correct validator proposes a block with a block $b'$ in its causal predecessors, 
then all correct validators can eventually \primitivesmall-deliver or decide $b'$ (see \Cref{lemma:totality}).
Therefore, a proposal from a correct validator will always eventually satisfy the external validity predicate.

\paragraph{Slotless blocks.}
Recall that \primitive may deliver slotless blocks without slot numbers.
Slotless blocks are not directly assigned to any slot.
However, validators receive a \eagercommit notification as soon as they become finalized. 
Moreover, when the block $b'$ that references $b$ as a causal predecessor becomes
committed, 
validators receive a \commit notification about $b$ together with $b'$ (using any deterministic order among them). 

\begin{algorithm}
  \caption{\system implementation with validity predicates $\mathcal{P}, \mathcal{Q}$.}
  \label{alg:bbca-log}
  \scriptsize

  \label{alg:sb-impl-pt1}
  \begin{algorithmic}[1]
    \Implements
      \State \primitivelog \textsc{\logsmall}, executed by process $p$.
    \EndImplements
    \Uses
      \State \primitive
      \State Validated Byzantine Total Order Broadcast, \textsc{vbtob}
    \EndUses
    
    \Parameters
    	\State $\numprocesses$
      	\Comment{Number of nodes}
      	\State $\numfaults$
      	\Comment{Maximum number of faulty nodes}
     	\State $N$
     	\Comment{Set of all $\numprocesses$ nodes}
      	\State $\Delta$ 
      	\Comment{Progress timeout period}
		\State $pid\gets 0$
		\Comment{\primitive identifier for node $p$}
		\State \textit{log}
		\Comment{Map of blocks per sot number}
		\State \textit{toFinalize}$\gets\{\}$
      	\Comment{Map of boolean values per slot number}
      	\State \textit{unreferrenced}$\gets\emptyset$
      	\Comment{Set of blocks previously not referenced by some block} 
      	\State \textit{firstUncommitted}$\gets 0$
		\Comment{First uncommitted log slot number}
		\State \textit{finalized}$\gets\emptyset$
		\Comment{Set of finalized blocks}	
		\State \textit{toFinalize}$\gets\emptyset$
		\Comment{Set of (seq no, blocks) tuples pending finalization}	
    \EndParameters
    
    \Struct{Block}
      \State sender\Comment{Validator who proposes the block}
      \State command \Comment{Block metadata}
      \State predecessors \Comment{Block references}
      \State payload \Comment{Transaction references}
    \EndStruct

	\UponEvent{\init()}
		\State $T_f.\operatorname{reset}(\Delta)$
		\Comment{timer with duration $\Delta$ for the first not uncommitted slot}
	\EndUponEvent
	
	\UponEvent{\propose(b)}
		\State $\operatorname{propose}(b,\bot)$
	\EndUponEvent
	
	\UponEvent{\bbcadeliver(\_, sn, b)}
		\State $\operatorname{finalize}(sn, b)$
		\If{$b\notin$\textit{unreferrenced}}
			\State\textit{unreferrenced}$\gets \{b\}\cup$\textit{unreferrenced}
		\EndIf
	\EndUponEvent
	
    \UponEvent{log[\operatorname{\textit{firstUncommitted}}] \neq None}
        \State $b\gets\log[\operatorname{\textit{firstUncommitted}}]$
        \If{$b\neq\bot$}
        	\ForAll{b'}{\operatorname{traverse}(b)}
				\State \textbf{trigger} \commit($b'$)
			\EndForAll
		\EndIf
        \State $T_f.\operatorname{reset}(\Delta)$
        \State \textit{firstUncommitted} $\gets$\textit{firstUncommitted}$+1$
   	\EndUponEvent
	
	\UponEvent{T_f.\operatorname{timeout}}
		\State\textbf{trigger} \bbcaabort(\textit{firstUncommitted})
		\State $T_f.\operatorname{reset}(\Delta)$
		\State \textit{firstUncommitted}$\gets$\textit{firstUncommitted}$+1$
		\While{log[\operatorname{\textit{firstNotTimedOut}}] \neq None}
			\State \textit{firstUncommitted}$\gets$\textit{firstUncommitted}$+1$
		\EndWhile
	\EndUponEvent
	
  \algstore{log}
  \end{algorithmic}
\end{algorithm}

\begin{algorithm}
    \scriptsize
    \ContinuedFloat
    \begin{algorithmic}[1]
    \algrestore{log}
    
	\UponEvent{\bbcaannounce(sn, b, C_R, \mathcal{Q})}
		\State $b'.\operatorname{payload}\gets$ payload()
		\State $\operatorname{propose}(b', [\finalize,sn, b, C_R])$
	\EndUponEvent
	
	\State\textbf{upon event} $\bbcadeliver(\_, \_, b)$ \textbf{such that} $b$.{command}$=[\finalize,sn', b', C_R]$	\textbf{do}
		\State\hskip\algorithmicindent\textbf{if} \textit{toFinalize} $[sn'][q]=\bot$ \textbf{then}
		\State\hskip\algorithmicindent\hskip\algorithmicindent\textit{toFinalize}$[sn'][q]\gets sn', b', C_R $
		\State\hskip \algorithmicindent\textbf{end if}
	\\
	
	\UponExists{(id,sn,b)}{|\{p\in N | \operatorname{\textit{toFinalize}}[sn][p]=(sn,b,C_R)\}|>2f}
		\State $\mathcal{C}\gets$ \textit{toFinalize}$[sn]$
		\State \textbf{if exists} $(sn,m,C_R)\in$ \textit{toFinalize}$[sn]$ \textbf{such that} $C_R\neq\emptyset$ \textbf{then}
		\State \hskip \algorithmicindent \textbf{trigger} $\bcpropose([sn, m,\mathcal{C}], \mathcal{Q})$
		\State{\textbf{else}}
		\State \hskip \algorithmicindent \textbf{trigger} $\bcpropose([sn,\bot,\mathcal{C}], \mathcal{Q})$
		\State{\textbf{end if}}
	\EndUponExists
	
	\UponEvent{\bcdecide(sn, b)}
		\If{$b=\bot$}
			\State \textit{log}$[sn]\gets b$	
			\State \textbf{return}
		\EndIf
		\State $\textit{toFinalize}\gets\{b\}\cup(b.\operatorname{predecessors}\setminus\textit{finalized})\cup\textit{toFinalize}$
	\EndUponEvent
	
	\UponExists{b\in\textit{toFinalize}}{b.\operatorname{predecessors}\in\textit{finalized}}
		\State \textbf{trigger} $\bbcadeliver(\_, sn, b)$
	\EndUponExists
	
		
	\UponEvent{\bbcacast(\_, sn, b)}
		\ForAll{b'}{b.\operatorname{predecessors}}
			\If{$b'\notin\textit{finalized}$}
				\State\textbf{trigger} $\operatorname{timeout}(sn,b)$ \textbf{ after } $\Delta$ 
			\EndIf
		\EndForAll
	\EndUponEvent
	
	\UponEvent{\operatorname{timeout}(sn,b)}
			\If{$b\notin\textit{finalized}$}
				\State\textbf{trigger} $\bbcaabort(sn)$ 
			\EndIf
	\EndUponEvent
	
	\Function{propose}{b, c}
		\State $id\gets p||pid$
		\State $pid\gets pid+1$
		\State $sn\gets$ next()
		\State $R\gets$ unreferrenced
		\State unreferrenced$\gets\emptyset$
		\State $b.\operatorname{sender}, b.\operatorname{command}, b.\operatorname{predecessors}
				\gets p, c, R$
		\State \textbf{trigger} \bbcacast($id, sn, b, \mathcal{P}$)
	\EndFunction
	
	\FunctionNoArgs{payload}{}
		\State returns the next batch of transaction references, possibly empty
	\EndFunctionNoArgs
	
	\FunctionNoArgs{next}{}
		\State returns the next $sn$ which $p$ is assigned
	\EndFunctionNoArgs
	 
	\Function{traverse}{p}
		\State traverses a graph with root $p$ and returns a list of its nodes in a deterministic order
	\EndFunction
	
	\Function{finalize}{sn, b}
		\State $\textit{finalized}\gets\{b\}\cup\textit{finalized}$
		\If{$sn\neq\emptyslot$}
			\State \textit{log}$[sn]\gets b$	
			\State \textbf{trigger} \eagercommit($b'$)
		\EndIf
	\EndFunction
	\caption{\system implementation (continues).}	
	\label{alg:sb-impl-pt2}
  \end{algorithmic}
\end{algorithm}


